\section*{अध्याय \ref{ch:b2:crustal-melting} --- भूपर्पटी का गलन: मिग्मेटाइट और ग्रेनाइट}

\begin{solution}{exo:b2:crustal-melting:1}
(क) प्रति \qty{100}{\gram}: $n(\ce{Al2O3}) = 17.03/101.96 = \qty{0.16703}{mol}$,
$n(\ce{CaO}) = 0.09183$, $n(\ce{Na2O}) = 0.06783$, $n(\ce{K2O}) = 0.03076$।
A/CNK $= 0.16703/0.19042 = 0.877$; A/NK $= 0.16703/0.09859 = 1.694$।
(ख) A/CNK $< 1 <$ A/NK: \omterm{def:b2:crustal-melting:asi}{मेटा-ऐलुमिनीय}। जो कैल्सियम फ़ेल्डस्पार नहीं ले सकते, वह
कैल्सियम सिलिकेटों में जाता है: हॉर्नब्लेंड (या क्लाइनोपाइरॉक्सीन, टाइटेनाइट)।
\end{solution}

\begin{solution}{exo:b2:crustal-melting:2}
(क) मस्कोवाइट और गार्नेट फ़ेल्डस्पारों से अधिक ऐलुमिनियम रखते हैं: \omterm{def:b2:crustal-melting:asi}{अति-ऐलुमिनीय}; यह ग्रेनाइट
S-प्रकार का हो सकता है। (ख) हॉर्नब्लेंड और टाइटेनाइट फ़ेल्डस्पारों से अधिक कैल्सियम रखते हैं:
\omterm{def:b2:crustal-melting:asi}{मेटा-ऐलुमिनीय}। (ग) एजिरीन और सोडियमी एम्फ़ीबोल फ़ेल्डस्पार के ऐलुमिनियम के बिना
सोडियम रखते हैं: \omterm{def:b2:crustal-melting:asi}{अतिक्षारीय}।
\end{solution}

\begin{solution}{exo:b2:crustal-melting:3}
एक किलोमीटर चट्टान का भार $2800\times9.81\times1000 = \qty{2.747e7}{Pa}$ है, अतः
$P = 0.2747\,z$ (\unit{kbar}, $z$ \unit{km} में)। सॉलिडस रेखा $T = 680 - 7.5(P -
2) = 695 - 2.060\,z$ है। $10 + 20z = 695 - 2.060z$ रखने पर $z = 685/22.060 =
\qty{31.1}{km}$, $T = \qty{631}{\celsius}$, $P = \qty{8.5}{kbar}$ मिलते हैं।
\end{solution}

\begin{solution}{exo:b2:crustal-melting:4}
(क) सफ़ेद परतें \omterm{def:b2:crustal-melting:migmatite}{ल्यूकोसोम} हैं, काली धारियाँ \omterm{def:b2:crustal-melting:migmatite}{मेलानोसोम}, सामान्य नाइस
मीसोसोम। (ख) अलग-अलग, सुस्पष्ट \omterm{def:b2:crustal-melting:migmatite}{ल्यूकोसोम}, \omterm{def:b2:crustal-melting:migmatite}{मेलानोसोम} और मीसोसोम: एक \omterm{def:b2:crustal-melting:migmatite}{मेटाटेक्साइट},
स्ट्रोमैटिक संरचना के साथ। (ग) चट्टान \omterm{def:b2:crustal-melting:migmatite}{डायाटेक्साइट} (श्लीरेनी और नेबुलाइटी) बन गई है:
गलन और आगे बढ़ा, गहरे खनिज भी पिघले, और गलित अपनी जगह पर रहा,
जिससे परतें मिट गईं।
\end{solution}

\begin{solution}{exo:b2:crustal-melting:5}
(क) $M = 39.10 + 3\times24.31 + 26.98 + 3\times28.09 + 12\times16.00 + 2\times1.008 =
\qty{417.30}{\gram\per\mole}$; $w = 18.02/417.30 = \qty{4.32}{\percent}$।
(ख) $F = 0.25\times0.0432/0.040 = 0.270$, यानी द्रव्यमान के अनुसार \qty{27}{\percent} गलित: यह
समान चट्टानों पर बायोटाइट निर्जलीकरण-गलन प्रयोगों में उत्पन्न आयतन के अनुसार 20 से \qty{30}{\percent}
के भीतर है।
\end{solution}

\begin{solution}{exo:b2:crustal-melting:6}
(क) मोलर द्रव्यमान ये हैं: पोटैशियम फ़ेल्डस्पार \qty{278.35}{\gram\per\mole}, एल्बाइट 262.24, एनॉर्थाइट 278.22, मस्कोवाइट
398.33। (ख) मात्राएँ: पोटैशियम फ़ेल्डस्पार $30/278.35 = \qty{0.10778}{mol}$, एल्बाइट $0.13347$,
एनॉर्थाइट $0.01797$, मस्कोवाइट $0.01255$। प्रति सूत्र-इकाई लिखने पर, पोटैशियम फ़ेल्डस्पार का रूप $\tfrac12$\ce{K2O.}$\tfrac12$\ce{Al2O3} है,
एल्बाइट $\tfrac12$\ce{Na2O.}$\tfrac12$\ce{Al2O3}, एनॉर्थाइट \ce{CaO.Al2O3}, मस्कोवाइट
$\tfrac12$\ce{K2O.}$\tfrac32$\ce{Al2O3}। $n(\ce{Al2O3}) = 0.05389 +
0.06673 + 0.01797 + 0.01883 = \qty{0.15742}{mol}$; $n(\ce{K2O}) = 0.05389 + 0.00628 =
0.06017$, $n(\ce{Na2O}) = 0.06673$, $n(\ce{CaO}) = 0.01797$, योग $0.14487$।
अतः A/CNK $= 0.15742/0.14487 = 1.087$ है। (ग) मस्कोवाइट के बिना तीनों फ़ेल्डस्पार
ठीक 1 देते हैं (\cref{prop:b2:crustal-melting:feldspar}); मस्कोवाइट, जिसमें प्रति सूत्र-इकाई पोटैशियम फ़ेल्डस्पार से एक \ce{Al2O3}
अधिक है, चट्टान को \omterm{def:b2:crustal-melting:asi}{अति-ऐलुमिनीय} बनाता है।
\end{solution}

\begin{solution}{exo:b2:crustal-melting:7}
बेसाल्ट: $900 + 2.5\times(1200 - 850) = \qty{1775}{\mega\joule\per\cubic\metre}$। भूपर्पटी:
$2.5\times(850 - 650) + 0.20\times900 = \qty{680}{\mega\joule\per\cubic\metre}$।
$V_c/V_b = 1775/680 = 2.6$। यह अनुमान उस ऊष्मा की उपेक्षा करता है जो चालन द्वारा ऊपर
और बगल में रिस जाती है, और ऊष्मा के फैलने में लगने वाले समय की; वह यह भी मानता है कि
बेसाल्ट की सारी ऊष्मा पिघलने को तैयार भूपर्पटी तक पहुँचती है: यह एक ऊपरी सीमा है।
\end{solution}

\begin{solution}{exo:b2:crustal-melting:8}
(क) $n(\ce{Al2O3}) = 14.15/101.96 = 0.13878$, $n(\ce{CaO}) = 0.01753$, $n(\ce{Na2O}) =
0.07809$, $n(\ce{K2O}) = 0.04416$; A/CNK $= 0.13878/0.13978 = 0.993$, A/NK $=
0.13878/0.12225 = 1.135$: \omterm{def:b2:crustal-melting:asi}{मेटा-ऐलुमिनीय}, बस किसी तरह। (ख) \qty{14.30}{\percent} \ce{Al2O3} के साथ,
$0.14025/0.13978 = 1.003$; \qty{14.00}{\percent} के साथ, $0.13731/0.13978 = 0.982$। सूचकांक
प्रमाणित अनिश्चितता के भीतर 1 के दोनों ओर है: ऑब्सीडियन
मेटा-ऐलुमिनीय--अति-ऐलुमिनीय सीमा पर स्थित है, और उसका वर्ग दृढ़ नहीं है।
\end{solution}

\begin{solution}{exo:b2:crustal-melting:9}
(क) माध्य $\bar P = \qty{5.0}{kbar}$, $\bar w = \qty{10.045}{\percent}$;
$\sum(P - \bar P)(w - \bar w) = 17.62$, $\sum(P - \bar P)^2 = 14.0$; ढलान
$17.62/14.0 = \qty{1.259}{\percent}$ प्रति किलोबार, अंतःखंड $10.045 - 1.259\times5.0 =
\qty{3.75}{\percent}$। ढलान उद्धृत 1.27 से
\qty{1}{\percent} के भीतर सहमत है। (ख) $3.75 + 1.259\times6 = \qty{11.30}{\percent}$। (ग) $10/0.1130 =
\qty{88.5}{\gram}$ गलित। (घ) रेखा केवल 3 और \qty{8}{kbar} के बीच,
\qty{800}{\celsius} पर, के आँकड़ों पर टिकी है; कुछ भी यह सुनिश्चित नहीं करता कि उच्चतर दाब पर विलेयता रैखिक बनी रहे,
जहाँ गलित और तरल अधिक समान हो जाते हैं।
\end{solution}

\begin{solution}{exo:b2:crustal-melting:10}
(क) रेखा \qty{50}{\celsius} में \qty{4}{kbar} चढ़ती है: $P = 6 + 0.08\,(T - 750)$;
\qty{770}{\celsius} पर $P = \qty{7.6}{kbar}$। (ख) 7.6 से \qty{6}{kbar} तक
तापमान रेखा के अनुदिश 770 से \qty{750}{\celsius} तक गिरता है; मुक्त ऊष्मा,
$2.5\times20 = \qty{50}{\mega\joule\per\cubic\metre}$, $50/900 = 0.056$ पिघलाती है, यानी आयतन के अनुसार चट्टान का लगभग
\qty{6}{\percent}। (ग) यही वह खिड़की है जो
हिमालय के \omterm{def:b2:crustal-melting:leucogranite}{ल्यूकोग्रेनाइटों} के लिए प्रस्तावित है, 6 से \qty{8}{kbar} और 750 से \qty{770}{\celsius}: अकेला विसंपीडन
गलित बनाता है, पर वक्र के हर बार पार होने पर केवल कुछ प्रतिशत।
\end{solution}

\begin{solution}{exo:b2:crustal-melting:11}
\omterm{def:b2:magma-emplacement:bodies}{सिल} के लिए $\ell^2/\kappa = (2000)^2/10^{-6} = \qty{4e12}{s} = \qty{0.13}{Myr}$;
अधःस्तरित परत के लिए $(10^4)^2/10^{-6} = \qty{1e14}{s} = \qty{3.2}{Myr}$। दोनों
दस मिलियन वर्षों की तुलना में छोटे हैं: एक घटना भूपर्पटी को इतने समय तक पिघला नहीं रख सकती, इसके लिए
बार-बार ऊष्मा की आपूर्ति चाहिए, बार-बार के अंतर्वेधनों से या किसी
मोटी हुई भूपर्पटी के धीमे तापन से।
\end{solution}

\begin{solution}{exo:b2:crustal-melting:12}
(क) $5.0x + 2.0(1 - x) = 2.75$ से $x = 0.75/3.0 = 0.25$ मिलता है: नियोसोम का एक-चौथाई
\omterm{def:b2:crustal-melting:migmatite}{ल्यूकोसोम} है। (ख) $0.25\times0.20 + 0.75\times4.0 = 0.05 + 3.00 = \qty{3.05}{\percent}$, जैसा
मीसोसोम में है: बंद-तंत्र व्याख्या टिकती है। (ग) जनक के \qty{100}{\gram} के लिए
\qty{25}{\gram} गलित और \qty{75}{\gram} अवशेष बने; यदि $m$ ग्राम
गलित निकल गया, तो $(25 - m)/(100 - m) = 0.10$, अतः $25 - m = 10 - 0.1m$ और $m = 15/0.9 =
\qty{16.7}{\gram}$: गलित का दो-तिहाई अनावरण से निकल गया।
\end{solution}

\begin{solution}{pb:b2:crustal-melting:1}
\noindent\textbf{भाग I।}

\textbf{1.} $39.10 + 3\times26.98 + 3\times28.09 + 12\times16.00 + 2\times1.008 =
\qty{398.33}{\gram\per\mole}$।

\textbf{2.} $18.02/398.33 = \qty{4.52}{\percent}$।

\textbf{3.} $200\times0.0452 = \qty{9.04}{\gram}$ जल प्रति किलोग्राम शिस्ट।

\textbf{4.} $9.97 + (13.82 - 9.97)\times2/3 = \qty{12.54}{\percent}$। न्यूनतम गलित
$9.04/0.1254 = \qty{72.1}{\gram}$ प्रति किलोग्राम है: $F \ge \qty{7.21}{\percent}$
(\cref{prop:b2:crustal-melting:water})।

\textbf{5.} $9.04/0.060 = \qty{151}{\gram}$: $F = \qty{15.1}{\percent}$।

\textbf{6.} उदाहरण के लिए: चट्टान में कम मस्कोवाइट हो सकता है; मस्कोवाइट का कुछ भाग सॉलिडस से नीचे ही,
क्वार्ट्ज़ के साथ, पोटैशियम फ़ेल्डस्पार और ऐलुमिनोसिलिकेट
में टूट चुका हो सकता है, और उसका जल गलन शुरू होने से पहले
तरल के रूप में निकल गया हो; गलित में \qty{6}{\percent} से अधिक जल हो सकता है; कुछ जल बायोटाइट में रह सकता है।

\noindent\textbf{भाग II।}

\textbf{7.} $680 - 7.5\times(7.8 - 2) = \qty{636.5}{\celsius}$; चट्टान का शिखर
उससे \qty{43.5}{\celsius} नीचे है और वह मुक्त जल के साथ भी नहीं पिघल सकती।

\textbf{8.} $605 + (19.2 - 15.3)/(27 - 15.3)\times65 = \qty{626.7}{\celsius}$।
$t = 0$ पर आधार, \qty{595}{\celsius} पर, \qty{32}{\celsius} नीचे है;
\qty{10}{Myr} पर, \qty{656}{\celsius} पर, वह ऊपर है।

\textbf{9.} आर्द्र सॉलिडस पर गलन को मुक्त जल चाहिए, और गर्म की गई कायांतरी चट्टान में
बहुत कम होता है: मुक्त जल पहले गलित में खप जाता है, जो फिर अभ्रकों के टूटने तक
बढ़ना बंद कर देता है।

\textbf{10.} 30 और \qty{40}{Myr} के बीच: $30 + (800 - 795)/(845 - 795)\times10 =
\qty{31.0}{Myr}$।

\textbf{11.} आधार 70 और \qty{80}{Myr} के बीच \qty{12.5}{kbar} से नीचे जाता है:
$70 + (13.7 - 12.5)/(13.7 - 12.4)\times10 = \qty{79.2}{Myr}$, तापमान
$914 - 14\times0.923 = \qty{901}{\celsius}$, एम्फ़ीबोल सीमा के भीतर: आधार पर कोई एम्फ़ीबोलाइट
पिघल रहा होता।

\textbf{12.} $90 + (11.0 - 10)/(11.0 - 9.6)\times10 = \qty{97.1}{Myr}$, तापमान $867 -
53\times0.714 = \qty{829}{\celsius}$। यह मस्कोवाइट सीमा से ऊपर है: मस्कोवाइट
टूट चुका है। यह बायोटाइट विघटन के 800 से \qty{850}{\celsius} के भीतर है
(कम दाब पर मापा गया): बायोटाइट, शिस्ट का दूसरा जलयोजित खनिज, अपनी
सीमा पर है।

\noindent\textbf{भाग III।}

\textbf{13.} $n(\ce{Al2O3}) = 15.0/101.96 = \qty{0.14712}{mol}$, $n(\ce{CaO}) =
0.55/56.08 = 0.00981$, $n(\ce{Na2O}) = 3.4/61.98 = 0.05486$, $n(\ce{K2O}) = 4.9/94.20 =
0.05202$।

\textbf{14.} A/CNK $= 0.14712/0.11669 = 1.261$।

\textbf{15.} A/NK $= 0.14712/0.10688 = 1.377$; A/CNK $> 1$: \omterm{def:b2:crustal-melting:asi}{अति-ऐलुमिनीय}।

\textbf{16.} \omterm{def:b2:crustal-melting:asi}{अति-ऐलुमिनीय}, एक शिस्ट से: \omterm{def:b2:crustal-melting:types}{S-प्रकार का ग्रेनाइट}। उसका अधिशेष ऐलुमिनियम
मस्कोवाइट और गार्नेट (या कॉर्डिएराइट) में रहता है।

\textbf{17.} \ce{Al2O3} के $0.14712 - 0.11669 = \qty{0.03043}{mol}$; मस्कोवाइट के रूप में
$0.03043\times398.33 = \qty{12.1}{\gram}$, \qty{100}{\gram} ग्रेनाइट में।

\textbf{18.} ऑब्सीडियन (बस किसी तरह) \omterm{def:b2:crustal-melting:asi}{मेटा-ऐलुमिनीय} है, मेटा-अवसादी गलित के अधिशेष ऐलुमिनियम के
बिना: उसका स्रोत संभवतः भिन्न है। पर एक नमूना बहुत कम सिद्ध करता है:
हर ग्रेनिटॉइड अनुक्रम में हल्के \omterm{def:b2:crustal-melting:asi}{अति-ऐलुमिनीय} सदस्य होते हैं, और अनुक्रम को उसकी पूरी
प्रवृत्ति से आँका जाता है।

\noindent\textbf{भाग IV।}

\textbf{19.} $10^{12}\,\unit{\cubic\metre}\times2650 = \qty{2.65e15}{kg}$।

\textbf{20.} ग्रेनाइट गलित में से उसका जल घटाने पर बचता है: $2.65\times10^{15}/0.94 =
\qty{2.82e15}{kg}$ गलित।

\textbf{21.} शिस्ट के द्रव्यमान का केवल $0.151\times0.70$ गलित के रूप में निकलता है:
$2.82\times10^{15}/(0.151\times0.70) = \qty{2.67e16}{kg}$।

\textbf{22.} $2.67\times10^{16}/2800 = \qty{9.5e12}{\cubic\metre}$, यानी
\qty{9.5e3}{\cubic\kilo\metre}, \omterm{def:b2:magma-emplacement:bodies}{प्लूटॉन} का लगभग दस गुना।

\textbf{23.} $9.5\times10^3/5 = \qty{1.9e3}{\square\kilo\metre}$, लगभग
\qty{44}{km} भुजा का एक वर्ग।

\textbf{24.} $F = \qty{7.21}{\percent}$ के साथ स्रोत $15.1/7.21 = 2.1$ गुना
बड़ा है, लगभग \qty{2.0e4}{\cubic\kilo\metre}। प्रश्न~5 के
\omterm{def:b2:mantle-melting:fraction}{गलन-अंश} के साथ \omterm{def:b2:magma-emplacement:bodies}{प्लूटॉन} बनाने के लिए पिघली स्रोत-भूपर्पटी का आयतन
\textbf{\qty[text-series-to-math=true]{9.5e3}{\cubic\kilo\metre}} है।
\end{solution}
